\chapter{Creating Sections}\label{Chap0_Sections}

\color{blue}In this chapter, we show how to create sections and subsections in a chapter. We will also show how to properly name and cite sections and subsections. 
\color{black}
\section{First Section Title}\label{Section1}
\color{blue} The first letter of each word that comprises the title of the section should be capitalized (\underline{F}irst \underline{S}ection \underline{T}itle). \color{black} 

\section[Section With Extremely Long and Lengthy Title That Does Not Fit On One Line And You Need More Than One Line]{Section With Extremely Long and Lengthy Title That \\ \vskip -3mm Does Not Fit On One Line And You Need More Than One Line} \label{Section1_2}

\color{blue} The first letter of each word that comprises the title of the section should be capitalized (see code for details). Write your title in [] the way you want it to show in your table of content. Write your title in \{\} the way you want it to show in your document. 


Next, we will show how to indent. The first paragraph is not indented; the following paragraphs are indented. You should skip indentation when you start a new idea. See below: 


\noindent This is how you get rid of the indentation.

And this is how you include it.


\color{blue} \noindent Now, let's see how to create subsections. Only the first word in the subsection title is capitalized. See below: \color{black} 
 
\subsection{Title of first subsection} \label{Subsection1_1}
\color{blue} The first letter of the first word is capitalized (\underline{T}itle of first subsection). Next, we show how to create a subsection title that does not fit on one line. 


Here is how to create subsections with long titles that do not fit on one line: \color{black} 

\subsection[Subsection with extremely long and lengthy title that does not fit on one line and you need more than one line]{Subsection with extremely long and lengthy title that does not fit \\ \vskip -11.75mm \hskip 0.93em  on one line and you need more than one line } \label{Subsection1_2}

\color{blue} Write your title in [] the way you want it to show in your table of content. Write your title in \{\} the way you want it to show in your document (see code for details). \color{black} 



\section{Cite Section and Subsection}\label{Section2}
\color{blue} Use Section with a capital `S' and Subsection with a capital `S' when you are citing a section and a subsection, respectively (even if you are citing in the middle of the sentence). Here are some examples with the correct syntax:

	\begin{singlespace}

	\begin{enumerate}
		\item According to Section \ref{Section2} ...
		\item According to Sections \ref{Section1} and \ref{Section2} ...
		\item According to Subsection \ref{Subsection1_2} ...
		\item According to Subsections \ref{Subsection1_1} and \ref{Subsection1_2} ...
	\end{enumerate}
	\color{black}
	
\end{singlespace}

\color{red}Here are some \textit{wrong} syntax that you should avoid using: 

\begin{singlespace}
	\begin{enumerate}
		\item According to section \ref{Section2}
		\item According to sec \ref{Section2}
		\item According to sec. \ref{Section2}
		\item According to sec. (\ref{Section2})
		\item According to subsection \ref{Subsection1_2}
		\item According to subsec \ref{Subsection1_2}
		\item According to subsec. \ref{Subsection1_2}
		\item According to subsec. (\ref{Subsection1_2})
	\end{enumerate}
	\color{black}
	
\end{singlespace}


	\color{black}
